% ============================================================================
% Section VII: Comparison with the Chilean daily schedule (CEN / PLEXOS)
% ============================================================================

\section{Comparison with the Chilean Daily Schedule}
\label{sec:cases}

The Chilean system operator, the Coordinador El\'ectrico Nacional (CEN),
computes the daily operation schedule of the national system (SEN) with
PLEXOS.  For each schedule, the input data and the solution (dispatch,
line flows, losses and marginal costs per bus and hour) are available.
This makes it possible to compare an independent model against the
official one on the full national network.

\subsection{Setup}

The comparison uses the CEN schedules of 14 days between October~2025
and May~2026,
plus additional days for specific checks.  Each case has a horizon of
seven days at hourly resolution, with about 250 buses and 315 lines,
hydro cascades, batteries, unit commitment and reserves.  The PLEXOS
inputs are converted with \texttt{plexos2gtopt} in a mode that follows
the PLEXOS conventions as closely as possible, and solved by \gtopt{}
with the monolithic method (MIP gap 0.5\%).  The first day
of each horizon, which is the day actually scheduled, is compared bus
by bus and hour by hour.  The comparison covers results only, not computing times.

The CEN has used two loss models: a piecewise-linear model (in 2021) and
a generalized, single-pass loss model in the recent schedules.  \gtopt{}
is run with two of the models of Section~\ref{sec:losses}:
\begin{itemize}
  \item \textbf{piecewise}: the traditional piecewise model with free
    losses ($\varepsilon_\ell = 0$);
  \item \textbf{tangent}: the tangent outer approximation with a loss
    price $\varepsilon_\ell = 1$~\$/MWh.
\end{itemize}

\subsection{Losses}

Table~\ref{tab:cen_losses} compares the losses of the LP with the
physical losses $R_\ell f_\ell^2/V_\ell^2$ evaluated at the LP flows,
over the seven-day horizon of the 14 cases.  With free piecewise losses,
the LP losses are 2.5 times the physical value in the median case.  The
excess comes from the fictitious-sink mechanism of
Section~\ref{sec:losses}: filling expensive segments first and
circulating flow in both directions when the sum of the two end-bus
duals is negative.  On the week of 7~April~2026, for example, the
piecewise model loses 91~GWh more than the physical losses, 65~GWh of
them in hours with $\pi_{a_\ell}+\pi_{b_\ell}<0$; the tangent model
loses 6~GWh more, and spills the surplus renewable energy instead.

\begin{table}[t]
  \centering
  \caption{LP losses versus physical losses over the horizon
    (median and range over 14 cases)}
  \label{tab:cen_losses}
  \begin{tabular}{@{}lcc@{}}
    \toprule
    & Piecewise, free & Tangent, $\varepsilon{=}1$ \\
    \midrule
    LP / physical losses & 2.50 (1.05--3.78) & 1.03 (1.00--1.19) \\
    Excess over physical & 152\% & 5\% \\
    \bottomrule
  \end{tabular}
\end{table}

The CEN's own solutions show the same trade-off.  In the two 2021
schedules solved with the piecewise model, the PLEXOS losses during the
hours with zero marginal cost are five times the secant approximation.
The fictitious losses amount to 8.7 and 16.1~GWh, i.e.\ 4.9\% and 8.6\%
of the daily demand, and the prices stay at exactly zero, never
negative.  In the recent schedules, solved with the generalized loss
model, the fictitious losses disappear but the surplus has no sink
left.  Marginal costs then become negative: in most of these schedules,
slightly negative across the whole system during solar hours, and on
one day down to $-132$~\$/MWh in a congested pocket of the northern
grid.  Under surplus, a loss model therefore faces a dilemma: either
fictitious losses with zero prices, or physical losses with negative
prices.  The tangent model with a small loss price chooses the latter,
consistently with the CEN's current practice.

\subsection{Marginal Costs}

The first comparison of marginal costs showed a moderate agreement
(median correlation 0.39 and mean absolute error 21~\$/MWh over all
buses, 11~\$/MWh at the reference bus, Polpaico~220~kV).  Its largest
deviations were deeply negative prices at a small load bus in the north,
down to $-891$~\$/MWh, where PLEXOS never goes below
$-1.3$~\$/MWh.  Inspecting the LP traced them to the conversion of 173
lines whose limits PLEXOS does not enforce: PLEXOS operates some of
them well above their rating (e.g.\ 190~MW on a 100~MW transformer),
whereas the converter had kept the rating as a hard limit.  The
congested chain created an isolated pocket with negative prices.
Converting these lines as PLEXOS treats them (free up to three times the
rating, penalized up to six times) removes the artifact
(Table~\ref{tab:cen_cmg}).

\begin{table}[t]
  \centering
  \caption{Hourly marginal costs of day~1, \gtopt{} vs.\ PLEXOS (CEN),
    piecewise losses: correlation and mean absolute error [\$/MWh] over
    all buses, with unenforced lines as hard limits (strict) and as
    PLEXOS treats them (extended)}
  \label{tab:cen_cmg}
  \begin{tabular}{@{}lcccc@{}}
    \toprule
    & \multicolumn{2}{c}{Strict} & \multicolumn{2}{c}{Extended} \\
    Day & Corr. & MAE & Corr. & MAE \\
    \midrule
    2026-04-07 & 0.33 & 20.9 & 0.69 & 11.3 \\
    2026-04-22 & 0.18 & 37.1 & 0.69 & 19.0 \\
    2026-05-25 & 0.32 & 45.9 & 0.49 & 43.3 \\
    \bottomrule
  \end{tabular}
\end{table}

With the tangent loss model and the extended lines, the correlation on
7~April reaches 0.72 and the mean absolute error 9.6~\$/MWh; the lowest
price in the north is $-8.1$~\$/MWh against $-13.9$~\$/MWh in PLEXOS.
The remaining differences come mainly from three sources:
\begin{itemize}
  \item \textbf{Constraints not yet converted.}  The $-132$~\$/MWh
    pocket of 25~May is not reproduced by any configuration (\gtopt{}
    gives 0 to $+9$~\$/MWh at those buses).  It is probably set by an
    operator-defined constraint or flowgate of the PLEXOS model that is
    not yet active in the conversion.
  \item \textbf{Conversion of model features.}  Several PLEXOS
    constructs (e.g.\ inflexible gas contracts, storage end conditions,
    reserve provisions of new batteries) have no one-to-one equivalent,
    and each conversion choice can move the marginal unit.
  \item \textbf{Degeneracy.}  With many units at equal cost and
    abundant zero-cost energy, the optimal dispatch and duals are not
    unique, and two exact solvers can legitimately report different
    marginal costs.
\end{itemize}

This comparison is ongoing work.  Its value is less in the agreement
figures than in making the differences explainable.  With an open model
whose formulation and LP can be inspected, each discrepancy can be
traced to a modeling choice in the converter, in \gtopt{} or in the
official model.  The artifact above, for instance, was found and
removed by reading the duals of a few saturated lines in the LP.
